% Canonical Paper I appendix.
The examples below separate several equality tests without developing a general
relation or loop theory.  In each case, histories are marked chronological
words, operators are affine maps, and endpoint values are obtained only after
choosing an input.

\subsection{Same value, different operators}

Let
\[
    \gamma=(\mathsf A_1),
    \qquad
    \delta=(\mathsf M_{\log 2}).
\]
At the input \(x=1\), both histories have endpoint value \(2\).  Their induced
operators are nevertheless
\[
    \nu_x(\gamma)=x+1,
    \qquad
    \nu_x(\delta)=2x,
\]
which are unequal.  Equality at one input therefore does not imply operator
equality.

\subsection{Same operator, different histories}

The marked histories
\[
    \gamma=(\mathsf A_1,\mathsf A_2),
    \qquad
    \delta=(\mathsf A_2,\mathsf A_1)
\]
are different chronological words, but both induce \(x\mapsto x+3\).  They
also have the same total charge.  Thus operator equality does not recover a
marked history, even when charge data are retained.

\subsection{Same charge, nonzero relative torsion}

For \(p,q\ne0\), the histories
\[
    \gamma=(\mathsf A_p,\mathsf M_q),
    \qquad
    \delta=(\mathsf M_q,\mathsf A_p)
\]
have the same ACS endpoint \((p,q)\).  Their operators differ by the translation
\[
    \tau(\gamma,\delta)=p(e^q-1).
\]
The ACS endpoint is therefore insensitive to order, whereas evaluation is not.
The weighted filling, rather than the terminal charge, detects the difference.

\subsection{Zero endpoint defect, nontrivial history}

For \(p\ne0\), the nonempty history
\[
    \ell=(\mathsf A_p,\mathsf A_{-p})
\]
induces the identity operator and hence has zero endpoint defect for every
input.  It is nevertheless not the empty marked history.  In particular,
operator neutrality does not imply literal or historical equality.

A slightly different phenomenon occurs for
\[
    (\mathsf A_p,\mathsf M_q,\mathsf A_{-p},\mathsf M_{-q}).
\]
This history has zero total ACS charge but induces translation by
\(p(1-e^{-q})\).  Thus charge neutrality does not imply operator neutrality.

\subsection{Open torsion and closed-loop drift}

For additive side \(p=\mu h\) and logarithmic multiplicative side
\(q=\lambda k\), the two open orders differ by
\[
    T_{\mathrm{open}}=p(e^q-1),
\]
whereas the chronological closed loop has vertical drift
\[
    T_{\mathrm{closed}}=p(1-e^{-q})=e^{-q}T_{\mathrm{open}}.
\]
They have the common first-order density \(\mu\lambda\), but they are distinct
finite quantities.  Consequently local order torsion, raw closed-loop
holonomy, and infinitesimal curvature should not be used as interchangeable
notions without the normalization or limiting operation that relates them.

These examples exhibit the hierarchy needed in the main text:
\[
    \text{marked history}
      \longrightarrow \text{operator}
      \longrightarrow \text{endpoint value},
\]
with total charge providing an additional, order-forgetting projection.  No
converse implication between adjacent levels holds in general.
